%%%%%%%%%%%%%%%%%%%%%%%%%%%%%%%%%%%%%%%%%%%%%%%%%%%%%%%%%%%%%%%%%%%%%%%%%%%%%%%
%  Laboratorio 2 - Filtros FIR con CMSIS-DSP en MCXN947
%  Estructura de carpetas esperada (igual que en el TP1):
%     main.tex
%     informe-fcefyn.sty      -> copiar de documentacion/TP1/
%     img/                    -> logos, capturas de fdatool, Config Tools y
%                                del osciloscopio
%
%  Imagenes que usa este documento (las que falten aparecen como un recuadro
%  rojo con el nombre esperado):
%     img/unc_logo.png, img/fcefyn_logo.png      (los mismos del TP1)
%     img/peripherals-lpadc-tp2.png              Config Tools > ADC0 en 16 bits
%     img/fda-pb-8k.png                          fdatool: respuesta del PB a 8 kS/s
%     img/fda-pa-8k.png                          fdatool: respuesta del PA a 8 kS/s
%     img/fda-pbanda-8k.png                      fdatool: respuesta del PBanda a 8 kS/s
%     img/fda-eb-8k.png                          fdatool: respuesta del EB a 8 kS/s
%     img/osc-bypass.jpeg                        entrada y salida en bypass
%     img/osc-pb.jpeg                            entrada y salida con PB
%     img/osc-pa.jpeg                            entrada y salida con PA
%     img/osc-pbanda.jpeg                        entrada y salida con PBanda
%     img/osc-eb.jpeg                            entrada y salida con EB
%     img/osc-tiempo-isr.jpeg                    pin de prueba: duracion de la ISR
%
%  Todo lo que aparece entre corchetes rojos (\completar{...}) es contenido
%  que falta llenar. Para la entrega, descomentar \entregafinal: el documento
%  avisara en el log si quedo alguna marca sin resolver.
%%%%%%%%%%%%%%%%%%%%%%%%%%%%%%%%%%%%%%%%%%%%%%%%%%%%%%%%%%%%%%%%%%%%%%%%%%%%%%%
\documentclass[11pt,a4paper]{article}
\usepackage{informe-fcefyn}

% \entregafinal

% ------------------------- Datos de la caratula ------------------------------
\materia{Procesamiento Digital de Señales}
\carrera{Ingeniería en Computación}
\subtitulo{Laboratorio 2}
\tema{Filtros FIR con CMSIS-DSP en MCXN947}
\plataforma{FRDM-MCXN947 \textperiodcentered{} MCUXpresso IDE \textperiodcentered{} CMSIS-DSP \textperiodcentered{} MATLAB}
\autor{Costamagna, Matias Javier}{44.512.437}{matias.costamagna@mi.unc.edu.ar}
\autor{Moreyra, Julián}{45.214.537}{julian.moreyra@mi.unc.edu.ar}
\autor{Rivera, Mariano}{45.432.921}{luismarianorivera.25@mi.unc.edu.ar}
\autor{Rojas, Facundo Nicolás}{44.629.222}{facundo.rojas.222@mi.unc.edu.ar}
\docentes{Rossi, Roberto Raúl \quad\textperiodcentered\quad Parlanti, Gustavo Fernando \quad\textperiodcentered\quad Molina, Germán Alcides \quad\textperiodcentered\quad Ferreyra, Pablo Alejandro}
\fechaentrega{\today}
\logoizquierdo{img/unc_logo.png}
\logoderecho{img/fcefyn_logo.png}

% ---------------- Parametros del diseno (usados en las tablas) ---------------
% Igual que en el TP1: CTIMER0 a 150 MHz, match 3 en toggle.
\frecuenciatimer{150000000}
\divisordisparo{2}

% ---------------- Carga de CPU estimada por el filtro FIR --------------------
% Reloj del nucleo [Hz] y ciclos por coeficiente de arm_fir_q15 con
% blockSize = 1. El 5 es la estimacion que usa generar_filtros.m; si se mide
% el valor real (seccion de verificacion), cambiarlo aca y la tabla de carga
% se recalcula sola.
\newcommand{\fcpu}{150000000}
\newcommand{\ciclosportap}{5}
% \carga{fs}{taps}: porcentaje del periodo de muestreo que ocupa el filtro.
% Se marca en rojo si supera el 100 %.
\newcommand{\carga}[2]{%
  \xdef\cargatmp{\fpeval{round(100*(#2)*\ciclosportap*(#1)/\fcpu,1)}}%
  \ifdim\cargatmp pt>100pt
    \textcolor{rojoaviso}{\textbf{\num{\cargatmp}}}%
  \else
    \num{\cargatmp}%
  \fi}
% \filacarga{fs}{taps PB}{taps PA}{taps PBanda}{taps EB}
\newcommand{\filacarga}[5]{%
  \xdef\ciclostmp{\fpeval{round(\fcpu/(#1))}}%
  \qty{#1}{\hertz} & \num{\ciclostmp} &
  \carga{#1}{#2} & \carga{#1}{#3} & \carga{#1}{#4} & \carga{#1}{#5}\\}
% \filataps{fs}{taps PB}{taps PA}{taps PBanda}{taps EB}
\newcommand{\filataps}[5]{%
  \xdef\tapstmp{\fpeval{#2+#3+#4+#5}}%
  \qty{#1}{\hertz} & \num{#2} & \num{#3} & \num{#4} & \num{#5} &
  \num{\tapstmp}\\}

% \filakaiser{filtro}{A [dB]}{transicion [Hz]}{N impl. a 8k}{N impl. a 48k}
% Orden estimado de la ventana de Kaiser: M = (A - 8) / (2,285 * 2 pi df / fs)
\newcommand{\filakaiser}[5]{%
  \xdef\kaiserocho{\fpeval{ceil(((#2)-8)/(2.285*2*pi*(#3)/8000))}}%
  \xdef\kaisercuarentayocho{\fpeval{ceil(((#2)-8)/(2.285*2*pi*(#3)/48000))}}%
  #1 & \qty{#3}{\hertz} & \qty{#2}{\deci\bel} &
  \num{\kaiserocho} & \num{#4} & \num{\kaisercuarentayocho} & \num{#5}\\}

\begin{document}

\caratula
\setcounter{tocdepth}{2}
\tableofcontents
\clearpage

%%%%%%%%%%%%%%%%%%%%%%%%%%%%%%%%%%%%%%%%%%%%%%%%%%%%%%%%%%%%%%%%%%%%%%%%%%%%%%%
\section{Introducción}
%%%%%%%%%%%%%%%%%%%%%%%%%%%%%%%%%%%%%%%%%%%%%%%%%%%%%%%%%%%%%%%%%%%%%%%%%%%%%%%

\begin{consigna}[Laboratorio 2 \textemdash{} Filtros FIR por muestras]
\textbf{Objetivo.} Realizar el diseño e implementación de filtros FIR usando
procesamiento por muestras.

\textbf{Descripción.} Realizar un programa aplicativo que sea capaz de aplicar
un filtro FIR, por muestras, a un buffer de memoria de 512 muestras, que es
adquirido con el Laboratorio~1 teniendo en cuenta las frecuencias de muestreo
de \qtylist{8;16;22;44;48}{\kilo\hertz}. Con una de las teclas de la placa de
evaluación FRDM-MCXN947 se habilitará la aplicación del filtro o se hará
\emph{bypass} del filtro a un buffer de salida distinto del buffer de entrada.
De este buffer de salida se enviarán las muestras al DAC de la placa.
Visualizar el resultado utilizando un osciloscopio y comparar los resultados.

\textbf{Requerimientos de los filtros.}
\begin{enumerate}[nosep,label=\textbf{F\arabic*},leftmargin=2.2em]
  \item \label{filtro:pb} Pasa bajos: $f_c = \qty{3600}{\hertz}$,
        $A_{\text{stop}} = \qty{30}{\deci\bel}$.
  \item \label{filtro:pa} Pasa altos: $f_c = \qty{35}{\hertz}$,
        $A_{\text{stop}} = \qty{30}{\deci\bel}$.
  \item \label{filtro:pbanda} Pasa banda: $f_{c1} = \qty{35}{\hertz}$,
        $f_{c2} = \qty{3500}{\hertz}$, $A_{\text{stop}} = \qty{30}{\deci\bel}$.
  \item \label{filtro:eb} Elimina banda: $f_r = \qty{50}{\hertz}$,
        $BW = \qty{15}{\hertz}$, $A_{\text{stop}} = \qty{25}{\deci\bel}$.
\end{enumerate}
\end{consigna}

\subsection{Requerimientos}

De la consigna se desprenden los siguientes requerimientos, que se usan como
guía para el diseño y como lista de verificación en la
sección~\ref{sec:verificacion}:

\begin{enumerate}[label=\textbf{R\arabic*},leftmargin=2.2em]
  \item \label{req:fs} Adquirir la señal con la aplicación del
        Laboratorio~1, a \qtylist{8;16;22;44;48}{\kilo\sps}, en un buffer de
        512 muestras.
  \item \label{req:diseno} Diseñar, para cada frecuencia de muestreo, filtros
        FIR que cumplan las especificaciones \ref{filtro:pb} a
        \ref{filtro:eb}.
  \item \label{req:cmsis} Aplicar el filtro por muestras al buffer de
        entrada, escribiendo el resultado en un buffer de salida distinto.
  \item \label{req:tecla-filtro} Con una tecla, habilitar el filtro o hacer
        \emph{bypass} del buffer de entrada al de salida.
  \item \label{req:dac} Enviar las muestras del buffer de salida al DAC.
  \item \label{req:osc} Visualizar la entrada y la salida con un osciloscopio
        y comparar los resultados.
\end{enumerate}

\begin{nota}[Extensión respecto de la consigna]
La consigna pide habilitar o deshabilitar \emph{un} filtro. La implementación
lo extiende: la tecla SW2 recorre en forma circular el \emph{bypass} y los
cuatro filtros \ref{filtro:pb}--\ref{filtro:eb}, de modo que los cuatro se
pueden comparar sobre la misma señal sin reprogramar la placa. Para el
filtrado se usa la biblioteca CMSIS-DSP en aritmética q15, con el mismo
formato del buffer del Laboratorio~1.
\end{nota}

\subsection{Correspondencia con la presentación}

\begin{table}[H]
\centering
\caption{Dónde se encuentra cada ítem pedido en la presentación.}
\label{tab:entregables}
\begin{tabular}{@{}l l@{}}
\toprule
\textbf{Ítem pedido} & \textbf{Ubicación en este informe}\\
\midrule
Configuración del ambiente          & Sección~\ref{sec:ambiente}\\
Diseño de los filtros               & Sección~\ref{sec:filtros}\\
Diagrama de diseño de la aplicación & Sección~\ref{sec:diseno}\\
Código fuente del proyecto          & Sección~\ref{sec:implementacion} y
                                      Anexo~\ref{sec:codigo}\\
Programa corriendo                  & Demostración en clase; evidencia en
                                      sección~\ref{sec:verificacion}\\
\bottomrule
\end{tabular}
\end{table}

%%%%%%%%%%%%%%%%%%%%%%%%%%%%%%%%%%%%%%%%%%%%%%%%%%%%%%%%%%%%%%%%%%%%%%%%%%%%%%%
\ejercicio{Configuración del ambiente}\label{sec:ambiente}
%%%%%%%%%%%%%%%%%%%%%%%%%%%%%%%%%%%%%%%%%%%%%%%%%%%%%%%%%%%%%%%%%%%%%%%%%%%%%%%

\subsection{Herramientas utilizadas}

\begin{table}[H]
\centering
\caption{Hardware y software del entorno de trabajo.}
\label{tab:herramientas}
\begin{tabular}{@{}l l@{}}
\toprule
\textbf{Elemento} & \textbf{Versión / modelo}\\
\midrule
Placa de evaluación     & FRDM-MCXN947 (Cortex-M33, núcleo \texttt{cm33\_core0})\\
IDE                     & MCUXpresso IDE v25.6.136\\
SDK                     & \texttt{SDK\_26.06.00\_FRDM-MCXN947}\\
Herramientas de config. & MCUXpresso Config Tools (Clocks, Pins, Peripherals)\\
Proyecto base           & Proyecto del Laboratorio 1 (\archivo{TP1})\\
Bibliotecas             & Drivers del SDK y CMSIS-DSP (\texttt{arm\_fir\_q15})\\
Diseño de filtros       & MATLAB R2023a (9.14.0.2206163), \texttt{fdatool} / \texttt{filterDesigner}\\
Generador de señales    & \\
Osciloscopio            & \\
\bottomrule
\end{tabular}
\end{table}

\subsection{Creación del proyecto}

El proyecto \archivo{TP2} parte del proyecto del Laboratorio 1: conserva la
cadena CTIMER0 $\rightarrow$ INPUTMUX $\rightarrow$ LPADC0 $\rightarrow$ DAC0,
la selección circular de la frecuencia de muestreo y la indicación con el LED.
Sobre esa base:

\begin{enumerate}[nosep]
  \item Se agregó al proyecto el código fuente de CMSIS-DSP (carpeta
        \archivo{CMSIS/DSP}) para disponer de \texttt{arm\_fir\_init\_q15()} y
        \texttt{arm\_fir\_q15()}.
  \item Los filtros se diseñaron en MATLAB y sus coeficientes se exportaron a
        un archivo de cabecera, \archivo{source/filtros.h}, con el script
        \archivo{generar_filtros.m} (sección~\ref{sec:filtros}).
  \item La lógica de la aplicación se escribió en \archivo{source/TP2.c}.
  \item La depuración se hace con el MCU-Link integrado de la placa, con la
        configuración guardada en \archivo{TP2 LinkServer Debug.launch}.
\end{enumerate}

\subsection{Relojes y pines}

Los relojes y la asignación de pines son los mismos del Laboratorio 1 (PLL0 a
\qty{150}{\mega\hertz} para el núcleo y el CTIMER0, FRO\_HF de
\qty{48}{\mega\hertz} para el ADC0 y el DAC0). La tabla~\ref{tab:pines} los
resume; las capturas de Config Tools están en el informe del Laboratorio 1.

\begin{table}[H]
\centering
\caption{Asignación de pines (sin cambios respecto del Laboratorio 1, salvo
la función de SW2).}
\label{tab:pines}
\small
\begin{tabular}{@{}l l l l@{}}
\toprule
\textbf{Función} & \textbf{Pin} & \textbf{Periférico} & \textbf{Conector / observación}\\
\midrule
Entrada analógica          & \pin{ADC0\_A0}          & LPADC0 canal A0    & \texttt{J4[2]} (Arduino A0), \qtyrange{0}{3.3}{\volt}\\
Salida analógica           & \pin{P4\_2}             & DAC0\_OUT          & \texttt{J1[4]}\\
Disparo del ADC            & \pin{P0\_19}            & CT0\_MAT3          & \texttt{J2[13]}, punto de prueba\\
LED rojo / verde / azul    & \pin{P0\_10}, \pin{P0\_27}, \pin{P1\_2} & GPIO0, GPIO1 & Activos en bajo\\
Tecla SW3 (velocidad)      & \pin{P0\_6}             & GPIO0, canal IRQ 0 & Flanco descendente, pull-up interno\\
Tecla SW2 (filtro)         & \pin{P0\_23}            & GPIO0, canal IRQ 1 & Flanco descendente\\
\bottomrule
\end{tabular}
\end{table}

\subsection{Cambios en los periféricos}\label{sec:cambios-perif}

La tabla~\ref{tab:cambios} resume lo que cambió respecto del Laboratorio 1.
El resto de la configuración (CTIMER0 en modo toggle sobre el match 3,
disparo del ADC por INPUTMUX, interrupciones de las teclas en las líneas 0
y 1 del GPIO0) se mantiene.

\begin{table}[H]
\centering
\caption{Diferencias de configuración respecto del Laboratorio 1.}
\label{tab:cambios}
\footnotesize
\begin{tabular}{@{}l p{5.0cm} p{6.4cm}@{}}
\toprule
\textbf{Elemento} & \textbf{Laboratorio 1} & \textbf{Laboratorio 2}\\
\midrule
Referencia del LPADC0 & \texttt{kLPADC\_ReferenceVoltageAlt1}
                      & \texttt{kLPADC\_ReferenceVoltageAlt3}\\
Calibración del ADC   & En \texttt{main()}
                      & En \texttt{ADC0\_init()} (generado por Config Tools)\\
Referencia del DAC0   & \texttt{kDAC\_ReferenceVoltageSourceAlt3}
                      & \texttt{kDAC\_ReferenceVoltageSourceAlt1}, como el
                        ejemplo \texttt{dac\_basic} del SDK\\
Inicialización del DAC0 & Config Tools
                      & \texttt{init\_dac()} antes de
                        \texttt{BOARD\_InitBootPeripherals()}\\
Tecla SW2             & Run/Stop & Selección de filtro (bypass, PB, PA, PBanda, EB)\\
Buffers               & \texttt{samples[512]}
                      & \texttt{buffer\_in[512]} y \texttt{buffer\_out[512]}\\
\bottomrule
\end{tabular}
\end{table}

El LPADC0 trabaja en resolución alta (\qty{16}{\bit}), así que su resultado,
centrado en cero, ocupa directamente todo el rango de un \texttt{q15\_t}
(sección~\ref{sec:q15}).

\captura{img/peripherals-lpadc-tp2.png}{Config Tools \textemdash{} ADC0}{Comando
de conversión del LPADC0 en resolución alta (\qty{16}{\bit}).}{fig:lpadc}

\subsubsection{Orden de inicialización}

\texttt{BOARD\_InitBootPeripherals()} arranca el CTIMER0 y habilita la
interrupción del ADC0, así que a partir de ese momento la rutina del ADC
empieza a ejecutarse y a escribir en el DAC. Por eso \texttt{main()} llama a
\texttt{init\_dac()} (reloj, alimentación desde el SPC e inicialización del
DAC0) \emph{antes} de inicializar los periféricos. Los filtros, en cambio, se
inicializan después: mientras el puntero al filtro activo sea \texttt{NULL},
la rutina del ADC hace \emph{bypass}.

\subsubsection{Prioridades de interrupción}

Igual que en el Laboratorio 1, las tres interrupciones (\texttt{ADC0\_IRQn},
\texttt{GPIO00\_IRQn} y \texttt{GPIO01\_IRQn}) quedan con la prioridad por
defecto (0), así que ninguna desaloja a otra. Esto es lo que permite que el
cambio de filtro se aplique de forma segura (sección~\ref{sec:cambio-filtro}).
A diferencia del Laboratorio 1, la rutina del ADC ahora puede ser larga
(ejecuta el filtro completo), y su duración frente al período de muestreo se
analiza en la sección~\ref{sec:carga}.

%%%%%%%%%%%%%%%%%%%%%%%%%%%%%%%%%%%%%%%%%%%%%%%%%%%%%%%%%%%%%%%%%%%%%%%%%%%%%%%
\ejercicio{Diseño de los filtros}\label{sec:filtros}
%%%%%%%%%%%%%%%%%%%%%%%%%%%%%%%%%%%%%%%%%%%%%%%%%%%%%%%%%%%%%%%%%%%%%%%%%%%%%%%

\subsection{Especificaciones}

Un filtro FIR de $N$ coeficientes calcula cada muestra de salida como
\begin{equation}
  y[n] = \sum_{k=0}^{N-1} b_k\, x[n-k].
  \label{eq:fir}
\end{equation}
Con coeficientes simétricos ($b_k = b_{N-1-k}$) la fase es lineal y todas las
frecuencias sufren el mismo retardo de grupo,
\begin{equation}
  \tau_g = \frac{N-1}{2\,f_s}.
  \label{eq:retardo}
\end{equation}

La consigna fija, para cada filtro, la frecuencia de corte (o la frecuencia
central y el ancho de banda, en el EB) y la atenuación mínima en la banda de
rechazo, pero no el ancho de la banda de transición ni el rizado admitido en
la banda de paso. Un filtro no queda definido sin esos dos datos, así que se
eligieron al diseñar en \texttt{fdatool}. La tabla~\ref{tab:specs} muestra
ambas partes. En todos los casos se fijó $A_{\text{pass}} = \qty{1}{\deci\bel}$
y se tomó $A_{\text{stop}}$ igual a la pedida.

\begin{table}[H]
\centering
\caption{Especificaciones de los filtros: lo que fija la consigna y los
bordes elegidos en \texttt{fdatool}. Las mismas bandas, en hertz, se usaron
para las cinco $f_s$.}
\label{tab:specs}
\small
\begin{tabular}{@{}l l l c c@{}}
\toprule
& \multicolumn{1}{c}{\textbf{Consigna}} &
  \multicolumn{3}{c}{\textbf{Diseño en \texttt{fdatool}}}\\
\cmidrule(lr){2-2}\cmidrule(l){3-5}
\textbf{Filtro} & \textbf{Bordes} &
\textbf{Bordes [Hz]} & \textbf{$A_{\text{pass}}$} & \textbf{$A_{\text{stop}}$}\\
\midrule
\ref{filtro:pb} PB         & $f_c = \qty{3600}{\hertz}$
                           & $f_{\text{pass}} = 3500$, $f_{\text{stop}} = 3700$
                           & \qty{1}{\deci\bel} & \qty{30}{\deci\bel}\\
\ref{filtro:pa} PA         & $f_c = \qty{35}{\hertz}$
                           & $f_{\text{stop}} = 20$, $f_{\text{pass}} = 35$
                           & \qty{1}{\deci\bel} & \qty{30}{\deci\bel}\\
\ref{filtro:pbanda} PBanda & \qtyrange{35}{3500}{\hertz}
                           & \begin{tabular}[c]{@{}l@{}}$f_{\text{stop1}} = 20$, $f_{\text{pass1}} = 35$\\
                             $f_{\text{pass2}} = 3500$, $f_{\text{stop2}} = 3700$\end{tabular}
                           & \qty{1}{\deci\bel} & \qty{30}{\deci\bel}\\
\ref{filtro:eb} EB         & \qty{50}{\hertz}, $BW = \qty{15}{\hertz}$
                           & \begin{tabular}[c]{@{}l@{}}$f_{\text{pass1}} = 35$, $f_{\text{stop1}} = 42{,}5$\\
                             $f_{\text{stop2}} = 57{,}5$, $f_{\text{pass2}} = 65$\end{tabular}
                           & \qty{1}{\deci\bel} & \qty{25}{\deci\bel}\\
\bottomrule
\end{tabular}
\end{table}

Los bordes se ubicaron así:
\begin{itemize}[nosep]
  \item En el PB, la frecuencia de corte de la consigna (\qty{3600}{\hertz})
        queda en el centro de la banda de transición de \qty{200}{\hertz}.
  \item En el PA y en el flanco inferior del PBanda, la banda de paso empieza
        exactamente en \qty{35}{\hertz} y la transición se ubicó por debajo, de
        \qtyrange{20}{35}{\hertz}. El flanco superior del PBanda usa la misma
        transición que el PB, de \qtyrange{3500}{3700}{\hertz}.
  \item En el EB, la banda de rechazo es exactamente la de la consigna,
        $\qty{50}{\hertz} \pm \qty{7.5}{\hertz}$, con transiciones de
        \qty{7.5}{\hertz} a cada lado.
\end{itemize}

\subsubsection{Método de diseño: ventana de Kaiser}

Los cuatro filtros se diseñaron por el método de ventaneo con ventana de
Kaiser: se trunca la respuesta impulsional ideal del filtro y se la multiplica
por una ventana cuyo parámetro $\beta$ regula el compromiso entre la
atenuación en la banda de rechazo y el ancho de la transición. Una
característica del método es que el rizado resulta igual en las dos bandas,
$\delta = \min(\delta_p, \delta_s)$, con
\begin{equation}
  \delta_p = \frac{10^{A_{\text{pass}}/20} - 1}{10^{A_{\text{pass}}/20} + 1}
           \approx 0{,}0575,
  \qquad
  \delta_s = 10^{-A_{\text{stop}}/20}.
\end{equation}
Con $A_{\text{stop}} = \qty{30}{\deci\bel}$, $\delta_s \approx 0{,}0316$ es
el más restrictivo, así que el rizado real de la banda de paso es de
\qty{0.55}{\deci\bel}, mejor que el \qty{1}{\deci\bel} pedido. En el EB
($A_{\text{stop}} = \qty{25}{\deci\bel}$, $\delta_s \approx 0{,}0562$) los dos
valores casi coinciden y el rizado de paso queda en \qty{0.98}{\deci\bel}. La
atenuación de diseño es $A = -20\log_{10}\delta$ y el parámetro de la ventana
resulta $\beta \approx 2{,}12$ para \qty{30}{\deci\bel} y
$\beta \approx 1{,}33$ para \qty{25}{\deci\bel}.

El orden $M$ que requiere la ventana de Kaiser se estima con
\begin{equation}
  M \approx \frac{A - 8}{2{,}285\;\Delta\omega},
  \qquad \Delta\omega = \frac{2\pi\,\Delta f}{f_s},
  \label{eq:kaiser}
\end{equation}
donde $\Delta f$ es la banda de transición más angosta del filtro. La
ecuación~\eqref{eq:kaiser} se usa en la sección~\ref{sec:taps} para explicar
la cantidad de coeficientes de cada filtro.

\paragraph{Comparación con el diseño equiripple.} La alternativa habitual,
y la que \texttt{fdatool} propone por defecto, es el diseño equiripple por el
algoritmo de Parks-McClellan. Ese método distribuye el error según la
tolerancia de cada banda en lugar de imponer el mismo rizado en ambas, así
que aprovecha la holgura de $A_{\text{pass}} = \qty{1}{\deci\bel}$ que la
ventana de Kaiser desperdicia. Su orden se puede estimar con la fórmula de
Kaiser para filtros óptimos,
\begin{equation}
  M_{\text{eq}} \approx
  \frac{-20\log_{10}\sqrt{\delta_p\,\delta_s} - 13}{14{,}6\;\Delta f/f_s},
  \label{eq:equiripple}
\end{equation}
que a \qty{8}{\kilo\sps} da unos 40, 530 y 870 coeficientes para el PB, el
PA (y el PBanda) y el EB, frente a 62, 818 y \num{1263} con la ventana de
Kaiser: entre un \qty{31}{\percent} y un \qty{36}{\percent} menos, en la
misma proporción para todas las $f_s$.

Aun así se eligió la ventana de Kaiser por su robustez con filtros muy
largos. El diseño por ventana es directo: se calcula la respuesta
impulsional ideal y se la multiplica por la ventana, sin ningún proceso
iterativo. Parks-McClellan, en cambio, es un algoritmo iterativo que, con
órdenes de varios miles como los que se necesitan a las $f_s$ altas (hasta
\num{7604} coeficientes en el EB a \qty{48}{\kilo\sps}), puede no converger o
dar resultados numéricamente imprecisos. El precio de esa robustez es entre
un \qty{30}{\percent} y un \qty{35}{\percent} más de coeficientes y, por lo
tanto, más carga de CPU. Sin embargo, con la estimación de carga de la
sección~\ref{sec:carga} ($c = 5$ ciclos por coeficiente), la diferencia sólo
cambia un caso: el EB a \qty{16}{\kilo\sps}, que pasaría de un
\qty{135}{\percent} a un \qty{93}{\percent} del período de muestreo. A partir
de \qty{22}{\kilo\sps}, el PA, el PBanda y el EB superan el período de
muestreo con cualquiera de los dos métodos.

Como el ADC puede muestrear a cinco velocidades, cada filtro se diseñó cinco
veces, una por $f_s$, con las mismas bandas expresadas en hertz. Así el filtro
seleccionado tiene la misma respuesta en frecuencia analógica cualquiera sea
la velocidad activa. El costo de esa decisión es que, al aumentar $f_s$, la
banda de transición ocupa una fracción menor del rango
$[0,\, f_s/2]$ y el orden necesario crece (sección~\ref{sec:taps}).

\captura[0.3\textheight]{img/fda-pb-8k.png}{fdatool}{Respuesta en magnitud
del filtro pasa bajos para $f_s = \qty{8}{\kilo\sps}$.}{fig:fda-pb}

\captura[0.3\textheight]{img/fda-pa-8k.png}{fdatool}{Respuesta en magnitud
del filtro pasa altos para $f_s = \qty{8}{\kilo\sps}$.}{fig:fda-pa}

\captura[0.3\textheight]{img/fda-pbanda-8k.png}{fdatool}{Respuesta en
magnitud del filtro pasa banda para $f_s = \qty{8}{\kilo\sps}$.}{fig:fda-pbanda}

\captura[0.3\textheight]{img/fda-eb-8k.png}{fdatool}{Respuesta en magnitud
del filtro elimina banda para $f_s = \qty{8}{\kilo\sps}$.}{fig:fda-eb}

\subsection{De MATLAB al firmware}

Los coeficientes de cada filtro se exportaron desde \texttt{fdatool} al
\emph{workspace} de MATLAB (\emph{File} $\rightarrow$ \emph{Export},
como \emph{Coefficients}) con el nombre \texttt{<tipo>\_<fs>}, por ejemplo
\texttt{PB\_8k} o \texttt{EB\_48k}. El script \archivo{generar_filtros.m}
toma las 20 variables y, para cada una:

\begin{enumerate}[nosep]
  \item Si algún coeficiente alcanza $|b_k| \ge 1$ (que no se puede
        representar en q15), escala todo el filtro para que el mayor quede en
        $32767/32768$. La variación de ganancia es menor a \qty{0.01}{\deci\bel}.
  \item Cuantiza a q15: $b_k^{q15} = \operatorname{round}\!\left(2^{15}\, b_k\right)$.
  \item Invierte el orden de los coeficientes, porque CMSIS-DSP los espera
        como $\{b_{N-1}, \dots, b_1, b_0\}$. En un filtro simétrico la
        inversión no cambia nada, pero el script no depende de eso.
  \item Si $N$ es impar agrega un coeficiente nulo, porque
        \texttt{arm\_fir\_init\_q15()} exige un número par de coeficientes
        mayor o igual a 4. Como el arreglo está invertido, ese cero queda
        como nuevo $b_0$: no modifica la respuesta en magnitud y sólo agrega
        una muestra de retardo.
  \item Escribe los arreglos \texttt{static const q15\_t} en el archivo de
        cabecera, junto con dos tablas indexadas por $[f_s][\text{tipo}]$:
        \texttt{coef\_tabla} (punteros a los coeficientes) y
        \texttt{taps\_tabla} (cantidad de coeficientes). También define
        \texttt{MAX\_TAPS}, el largo del filtro más largo.
\end{enumerate}

El script además imprime una estimación de la carga de CPU de cada filtro,
que se analiza en la sección~\ref{sec:carga}. El archivo generado
(\archivo{FiltrosTP2.h}) se incorporó al proyecto como
\archivo{source/filtros.h}.

\subsection{Cantidad de coeficientes y memoria}\label{sec:taps}

\begin{table}[H]
\centering
\caption{Cantidad de coeficientes de cada filtro (después del relleno a un
número par), tomada de \texttt{taps\_tabla} en \texttt{source/filtros.h}.}
\label{tab:taps}
\begin{tabular}{@{}l r r r r r@{}}
\toprule
\textbf{$f_s$} & \textbf{PB} & \textbf{PA} & \textbf{PBanda} & \textbf{EB} & \textbf{Total}\\
\midrule
\filataps{8000} {64} {822} {822} {1270}
\filataps{16000}{124}{1642}{1640}{2536}
\filataps{22000}{170}{2256}{2254}{3486}
\filataps{44000}{340}{4508}{4508}{6970}
\filataps{48000}{370}{4918}{4916}{7604}
\midrule
\multicolumn{5}{@{}l}{\textbf{Total}} & \num{51220}\\
\bottomrule
\end{tabular}
\end{table}

Se observan dos cosas:

\begin{itemize}
  \item La cantidad de coeficientes crece casi en proporción a $f_s$ (por
        ejemplo, el EB pasa de \num{1270} coeficientes a \qty{8}{\kilo\sps} a
        \num{7604} a \qty{48}{\kilo\sps}, unas 6 veces, igual que la relación
        entre las frecuencias de muestreo). Es lo esperable al mantener fijas
        las bandas en hertz: el ancho de transición normalizado
        $\Delta f / f_s$ se achica en la misma proporción y el orden de un FIR
        es inversamente proporcional a ese ancho.
  \item El PB es mucho más corto que los otros tres filtros, aunque su
        atenuación es la misma que la del PA y el PBanda. La diferencia está
        en el ancho de la banda de transición, como se explica a
        continuación.
\end{itemize}

La ecuación~\eqref{eq:kaiser} explica esa diferencia. La transición más
angosta de cada filtro (tabla~\ref{tab:specs}) es de \qty{200}{\hertz} en el
PB, \qty{15}{\hertz} en el PA y el PBanda (de \qtyrange{20}{35}{\hertz}) y
\qty{7.5}{\hertz} en el EB. La tabla~\ref{tab:kaiser} compara el orden que
predice la fórmula con la cantidad de coeficientes implementada, que es
$M + 1$ más el eventual cero de relleno.

\begin{table}[H]
\centering
\caption{Orden estimado con la ecuación~\eqref{eq:kaiser} y cantidad de
coeficientes implementada, a la menor y a la mayor $f_s$.}
\label{tab:kaiser}
\begin{tabular}{@{}l r r r r r r@{}}
\toprule
& & & \multicolumn{2}{c}{\textbf{\qty{8}{\kilo\sps}}} &
      \multicolumn{2}{c}{\textbf{\qty{48}{\kilo\sps}}}\\
\cmidrule(lr){4-5}\cmidrule(l){6-7}
\textbf{Filtro} & \textbf{$\Delta f$} & \textbf{$A$} &
\textbf{$M$ estimado} & \textbf{$N$ impl.} &
\textbf{$M$ estimado} & \textbf{$N$ impl.}\\
\midrule
\filakaiser{PB}    {30}{200}{64}  {370}
\filakaiser{PA}    {30}{15} {822} {4918}
\filakaiser{PBanda}{30}{15} {822} {4916}
\filakaiser{EB}    {25}{7.5}{1270}{7604}
\bottomrule
\end{tabular}
\end{table}

La estimación coincide con lo implementado con un error menor al
\qty{1}{\percent} en todos los casos, lo que confirma que los filtros del
firmware son los de la tabla~\ref{tab:specs} y que se usaron las mismas
bandas en hertz para todas las $f_s$.

La explicación de fondo está en dónde caen los bordes. El PB tiene su
transición en \qtyrange{3500}{3700}{\hertz} y puede ser de
\qty{200}{\hertz}. En cambio, el PA y el flanco inferior del PBanda pasan de
rechazo a paso entre \qty{20}{\hertz} y \qty{35}{\hertz}, y el EB tiene que
entrar y salir de una banda de rechazo de \qty{15}{\hertz} con transiciones
de \qty{7.5}{\hertz}. Según la ecuación~\eqref{eq:kaiser}, $M$ es
proporcional a $f_s/\Delta f$: una transición de unos pocos hertz con $f_s$
de miles de hertz requiere miles de coeficientes. Ese es el origen de la
carga de CPU que se analiza en la sección~\ref{sec:carga}.

A \qty{8}{\kilo\sps} el PB tiene además una particularidad: la frecuencia de
Nyquist es \qty{4000}{\hertz} y la banda de rechazo empieza en
\qty{3700}{\hertz}, así que sólo atenúa los últimos \qty{300}{\hertz} del
espectro representable. A esa $f_s$, su efecto sobre una señal de prueba es
pequeño y se observa recién por encima de \qty{3.5}{\kilo\hertz}.

Los \num{51220} coeficientes ocupan \num{102440} bytes (unos 100~KiB) de memoria flash, una fracción chica de los 2~MB del
MCXN947. En RAM sólo hace falta un buffer de estado de
$\texttt{MAX\_TAPS} + 1 = \num{7605}$ muestras (\num{15210} bytes), compartido
por todos los filtros (sección~\ref{sec:instancias}).

Con la ecuación~\eqref{eq:retardo}, el retardo de grupo resulta casi igual
para todas las $f_s$, porque $N$ crece en la misma proporción que $f_s$. El
filtro más largo, el EB, retrasa la señal unos
$\num{1269}/(2\cdot\qty{8}{\kilo\hertz}) \approx \qty{79}{\milli\second}$,
lo que se puede observar con el osciloscopio
(sección~\ref{sec:verificacion}).

%%%%%%%%%%%%%%%%%%%%%%%%%%%%%%%%%%%%%%%%%%%%%%%%%%%%%%%%%%%%%%%%%%%%%%%%%%%%%%%
\ejercicio{Diseño de la aplicación}\label{sec:diseno}
%%%%%%%%%%%%%%%%%%%%%%%%%%%%%%%%%%%%%%%%%%%%%%%%%%%%%%%%%%%%%%%%%%%%%%%%%%%%%%%

\subsection{Arquitectura general}

La figura~\ref{fig:bloques} muestra las partes de la aplicación. El camino de
la señal es el del Laboratorio 1 con un bloque de procesamiento en el medio:
el CTIMER0 dispara cada conversión a través del INPUTMUX, el LPADC0 deja el
resultado en su FIFO y la rutina de interrupción lo convierte a q15, lo pasa
por el filtro activo (o lo copia, en \emph{bypass}) y escribe el resultado en
el DAC0. Las teclas sólo modifican el estado: SW3 cambia la frecuencia de
muestreo (match del CTIMER0 y color del LED) y SW2 el tipo de filtro. Ambas
avisan a la rutina del ADC, que es la única que cambia el filtro activo. El
lazo principal queda vacío.

\begin{figure}[H]
\centering
\begin{tikzpicture}[x=1cm,y=1cm, every node/.append style={font=\footnotesize}]
  \tikzset{bloque/.append style={minimum width=1.9cm}}
  % --- Fila 1: camino de la señal ---
  \node[externo] (gen) at (0,0)     {Generador\\de señal};
  \node[hw]      (adc) at (3.2,0)   {\periferico{LPADC0}\\canal A0, \qty{16}{\bit}};
  \node[sw]      (isr) at (6.6,0)   {ISR\\del ADC0};
  \node[hw]      (dac) at (10.2,0)  {\periferico{DAC0}\\\qty{12}{\bit}};
  \node[externo] (osc) at (13.5,0)  {Osciloscopio};

  % --- Fila 2 ---
  \node[hw]  (mux)  at (3.2,-2.0)  {\periferico{INPUTMUX}};
  \node[sw]  (fir)  at (6.6,-2.0)  {\texttt{arm\_fir\_q15}\\o \emph{bypass}};
  \node[mem] (coef) at (10.2,-2.0) {\texttt{coef\_tabla}\\$[f_s][\text{tipo}]$, flash};
  \node[mem] (buf)  at (13.5,-2.0) {\texttt{buffer\_in}\\\texttt{buffer\_out}\\512 $\times$ q15};

  % --- Fila 3 ---
  \node[externo] (p19) at (0,-4.0)    {\pin{P0\_19}\\punto de prueba};
  \node[hw]      (tmr) at (3.2,-4.0)  {\periferico{CTIMER0}\\match 3};
  \node[sw]      (isk) at (6.6,-4.0)  {ISR de\\teclas};
  \node[hw]      (led) at (10.2,-4.0) {LED RGB};

  % --- Fila 4 ---
  \node[hw] (sw) at (6.6,-5.8) {Teclas\\SW3 (velocidad) \,/\, SW2 (filtro)};

  % --- Señal ---
  \draw[senal] (gen) -- node[etiqueta,above]{$x(t)$} (adc);
  \draw[senal] (adc) -- node[etiqueta,above]{FIFO0} (isr);
  \draw[senal] (isr) -- node[etiqueta,above]{código} (dac);
  \draw[senal] (dac) -- node[etiqueta,above]{$y(t)$} (osc);

  % --- Datos y disparo ---
  \draw[flecha,{Stealth[length=2mm]}-{Stealth[length=2mm]}]
       (isr) -- node[etiqueta,right]{$x[n]$, $y[n]$} (fir);
  \draw[flecha] (coef) -- node[etiqueta,above]{$b_k$} (fir);
  \draw[flecha] ([yshift=-3mm]isr.east) -- node[etiqueta,above,sloped,pos=0.75]{guarda} (buf.north west);
  \draw[flecha] (tmr) -- node[etiqueta,right]{MAT3 (toggle)} (mux);
  \draw[flecha] (mux) -- node[etiqueta,right]{trigger 0} (adc);
  \draw[flecha] (tmr) -- (p19);

  % --- Control ---
  \draw[control] (sw)  -- node[etiqueta,right]{flanco $\downarrow$} (isk);
  \draw[control] (isk) -- node[etiqueta,above]{nuevo MR} (tmr);
  \draw[control] (isk) -- node[etiqueta,above]{color} (led);
  \draw[control] (isk) -- node[etiqueta,right]{\texttt{cambio\_filtro}} (fir);

  % --- Leyenda ---
  \begin{scope}[shift={(1.6,-7.2)},every node/.style={font=\scriptsize}]
    \foreach \c/\t [count=\k from 0] in {bloquehw/Periférico,bloquesw/Firmware,
                                          bloquemem/Memoria}
      {\fill[\c,draw=azulunc] (2.4*\k,0) rectangle ++(0.35,0.25);
       \node[right] at (2.4*\k+0.4,0.125) {\t};}
    \draw[grismedio,dashed] (7.2,0) rectangle ++(0.35,0.25);
    \node[right] at (7.6,0.125) {Externo};
  \end{scope}
\end{tikzpicture}
\caption{Diagrama de bloques de la aplicación. Trazo continuo azul: señal;
negro: datos y disparo; punteado: control.}
\label{fig:bloques}
\end{figure}

\subsection{Velocidades de muestreo}

La generación de las velocidades no cambió respecto del Laboratorio 1: con el
timer contando a \ftimer{}, reiniciándose en el match y con la salida en
toggle,
\begin{equation}
  f_s = \frac{f_{\text{timer}}}{2\,(\text{MR} + 1)}.
  \label{eq:fs}
\end{equation}
La tabla~\ref{tab:fs} repite los valores de MR (arreglo \texttt{sample\_vel})
y el color del LED de cada estado.

\begin{table}[H]
\centering
\caption{Velocidades de muestreo, valor de match programado y color del LED.}
\label{tab:fs}
\begin{tablamuestreo}
  \filafs{8000} {9374}{Verde}   {0}{1}{0}
  \filafs{16000}{4687}{Amarillo}{1}{1}{0}
  \filafs{22000}{3408}{Azul}    {0}{0}{1}
  \filafs{44000}{1704}{Rojo}    {1}{0}{0}
  \filafs{48000}{1562}{Blanco}  {1}{1}{1}
\end{tablamuestreo}
\end{table}

\subsection{Máquina de estados}

El comportamiento de las teclas se modela con dos variables de estado
independientes (figura~\ref{fig:estados}):

\begin{itemize}[nosep]
  \item \texttt{state\_counter}: índice de velocidad (E0 a E4). Avanza en
        forma circular con SW3. Cada cambio reprograma el match del CTIMER0,
        actualiza el color del LED y obliga a cambiar los coeficientes del
        filtro activo por los diseñados para la nueva $f_s$.
  \item \texttt{filter\_counter}: tipo de filtro (F0 a F4). Avanza en forma
        circular con SW2: \emph{bypass} $\rightarrow$ PB $\rightarrow$ PA
        $\rightarrow$ PBanda $\rightarrow$ EB $\rightarrow$ \emph{bypass}.
\end{itemize}

Las dos teclas actúan en cualquier estado, y la combinación de ambas
variables elige uno de los 20 filtros de la tabla~\ref{tab:taps}. El LED sólo
indica la velocidad, no el filtro. Al reiniciar la placa el sistema arranca
en E0 (\qty{8}{\kilo\sps}) y en \emph{bypass}.

\begin{figure}[H]
\centering
\begin{tikzpicture}
  % --- Anillo de velocidades ---
  \foreach \i/\f/\r/\g/\b in {0/8/0/1/0, 1/16/1/1/0, 2/22/0/0/1,
                                3/44/1/0/0, 4/48/1/1/1} {
    \node[estado] (e\i) at ({90-72*\i}:2.5)
      {\textbf{E\i}\\[-1pt]{\footnotesize\f\,kS/s}\\[1pt]\led{\r}{\g}{\b}};
  }
  \foreach \i/\j in {0/1,1/2,2/3,3/4,4/0}
    \draw[flecha] (e\i) to[bend left=18] node[etiqueta,sloped,above]{SW3} (e\j);
  \node[etiqueta] at (0,0) {avance\\circular};
  \draw[flecha] (-2.2,3.3) node[left,etiqueta]{reset} -- (e0);

  % --- Anillo de filtros ---
  \begin{scope}[xshift=7.6cm]
    \foreach \i/\n in {0/Bypass, 1/PB, 2/PA, 3/PBanda, 4/EB} {
      \node[estado] (f\i) at ({90-72*\i}:2.5)
        {\textbf{F\i}\\[1pt]{\footnotesize\n}};
    }
    \foreach \i/\j in {0/1,1/2,2/3,3/4,4/0}
      \draw[flecha] (f\i) to[bend left=18] node[etiqueta,sloped,above]{SW2} (f\j);
    \node[etiqueta] at (0,0) {avance\\circular};
    \draw[flecha] (-2.2,3.3) node[left,etiqueta]{reset} -- (f0);
  \end{scope}
\end{tikzpicture}
\caption{Izquierda: selección circular de la velocidad de muestreo con SW3 y
color del LED en cada estado. Derecha: selección circular del filtro con SW2.}
\label{fig:estados}
\end{figure}

\subsection{Formato q15 y conversiones}\label{sec:q15}

Un número q15 es un entero de 16 bits con signo que representa el valor
\begin{equation}
  v_{\mathrm{q15}} = \frac{x_{q15}}{2^{15}} \in \left[-1,\; 1-2^{-15}\right].
\end{equation}
En resolución alta el LPADC entrega un resultado de \qty{16}{\bit} sin signo,
$v \in [0,\, 65535]$. Como la entrada es single-ended, $v$ es unipolar; para
centrarla en cero basta restar la mitad de escala:
\begin{equation}
  x[n] = v - 32768 \in [-32768,\; 32767].
\end{equation}
El resultado ocupa exactamente el rango de un \texttt{q15\_t}:
\qty{0}{\volt} corresponde a $-1$, \qty{1.65}{\volt} a $0$ y
\qty{3.3}{\volt} a casi $+1$.

A la salida, la muestra filtrada $y[n]$ (en q15) se lleva a los
\qty{12}{\bit} del DAC descartando los 4 bits menos significativos y
volviendo a sumar la mitad de escala, con saturación:
\begin{equation}
  \text{código}_{\text{DAC}} =
  \mathrm{sat}_{[0,\,4095]}\!\left(\frac{y[n]}{16} + 2048\right).
\end{equation}
La saturación es necesaria porque, a diferencia del Laboratorio 1, la salida
ya no es una copia de la entrada: un filtro con ganancia mayor a 1 en alguna
frecuencia, o el rizado de la banda de paso, puede llevar $y[n]$ fuera del
rango.

\subsection{Filtrado con CMSIS-DSP}

\subsubsection{Procesamiento muestra a muestra}

\texttt{arm\_fir\_q15()} implementa la ecuación~\eqref{eq:fir} sobre un
bloque de \texttt{blockSize} muestras. Con \texttt{BLOCK\_SIZE = 1} se llama
una vez por interrupción del ADC con una sola muestra, así que la salida se
obtiene en la misma interrupción en que llega la entrada y el único retardo
agregado es el del propio filtro (ecuación~\eqref{eq:retardo}). La función
guarda internamente las últimas $N$ muestras en el buffer de estado, que
cumple el papel de línea de retardo.

Internamente los coeficientes y las muestras son 1.15; cada producto es un
2.30 y se acumula en 64 bits (formato 34.30), de modo que la suma no desborda
aunque el filtro tenga miles de coeficientes. Al final el acumulador se
trunca a 15 bits fraccionarios y se satura a 1.15. La única pérdida de
precisión, además de la cuantización de los coeficientes, ocurre en ese
último paso.

\subsubsection{Instancias de los filtros}\label{sec:instancias}

Cada filtro de CMSIS-DSP se representa con una estructura
\texttt{arm\_fir\_instance\_q15} (cantidad de coeficientes, puntero a los
coeficientes y puntero al buffer de estado). Al arrancar,
\texttt{init\_filtros()} inicializa una matriz de $5\times4$ instancias,
\texttt{arreglo\_filtros[fs][tipo]}, a partir de \texttt{coef\_tabla} y
\texttt{taps\_tabla}. El filtro activo es el apuntado por
\texttt{pte\_aux}, de modo que cambiar de filtro es sólo cambiar un puntero.

Como hay un único filtro activo a la vez, las 20 instancias comparten el
mismo buffer de estado, \texttt{estado\_fir}, dimensionado para el filtro más
largo ($\texttt{MAX\_TAPS} + \texttt{BLOCK\_SIZE}$ muestras). Tener un buffer
por filtro costaría unos 100~KiB de RAM sin ningún beneficio, porque
el estado de un filtro inactivo no se usa.

\subsubsection{Cambio sincronizado de filtro}\label{sec:cambio-filtro}

Las rutinas de las teclas no tocan el filtro directamente: sólo actualizan
su variable de estado y ponen en 1 la bandera \texttt{cambio\_filtro}. La
rutina del ADC revisa la bandera antes de procesar cada muestra y, si está
activa, llama a \texttt{aplicar\_seleccion()}, que:

\begin{enumerate}[nosep]
  \item elige la instancia \texttt{arreglo\_filtros[state\_counter][filter\_counter-1]}
        (o ninguna, en \emph{bypass});
  \item borra el buffer de estado, para que el nuevo filtro no arranque con
        la historia del anterior (que puede haber sido calculada a otra $f_s$
        o con otra cantidad de coeficientes).
\end{enumerate}

Así el cambio ocurre siempre entre dos muestras y nunca en medio de un
\texttt{arm\_fir\_q15()}. Como todas las interrupciones tienen la misma
prioridad, tampoco puede pasar que una tecla modifique el estado mientras la
rutina del ADC lo está leyendo. El borrado del estado produce un transitorio
de $N$ muestras a la salida después de cada cambio, durante el cual la línea
de retardo se vuelve a llenar.

\subsection{Rutinas de interrupción}

La figura~\ref{fig:flujo-adc} muestra la rutina del ADC, que se ejecuta una
vez por muestra, y la figura~\ref{fig:flujo-teclas} las de las teclas.

\begin{figure}[H]
\centering
\begin{tikzpicture}[node distance=0.55cm and 0.4cm]
  \node[inicio]  (a0) {\texttt{ADC0\_IRQHandler}};
  \node[proceso, below=of a0] (a1) {Leer y limpiar banderas\\de trigger y de estado};
  \node[decision, below=of a1] (a2) {¿Hay resultado\\en la FIFO0?};
  \node[decision, below=0.6cm of a2] (a3) {¿\texttt{cambio\_filtro}?};
  \node[proceso, right=1.0cm of a3] (a4)
     {\texttt{aplicar\_seleccion()}\\elige instancia y\\borra el estado};
  \node[proceso, below=0.6cm of a3] (a5)
     {$x[n] \leftarrow \mathrm{res} - 32768$\\\texttt{buffer\_in[i]} $\leftarrow x[n]$};
  \node[decision, below=0.6cm of a5] (a6) {¿\emph{Bypass}?};
  \node[proceso, minimum width=3.0cm, below=0.7cm of a6, xshift=-2.3cm] (a7)
     {$y[n] \leftarrow x[n]$};
  \node[proceso, minimum width=3.0cm, below=0.7cm of a6, xshift=2.3cm] (a8)
     {\texttt{arm\_fir\_q15}\\(\texttt{pte\_aux}, 1 muestra)};
  \node[proceso, minimum width=6.4cm, below=2.0cm of a6] (a9)
     {\texttt{buffer\_out[i]} $\leftarrow y[n]$\\
      \texttt{DAC\_SetData}$\big(\mathrm{sat}(y[n]/16+2048)\big)$};
  \node[proceso, below=of a9] (a10) {$i \leftarrow (i+1) \bmod 512$};
  \node[inicio,  below=of a10] (a11) {Retorno};

  \draw[flecha] (a0) -- (a1);
  \draw[flecha] (a1) -- (a2);
  \draw[flecha] (a2) -- node[etiqueta,right]{sí} (a3);
  \draw[flecha] (a3) -- node[etiqueta,above]{sí} (a4);
  \draw[flecha] (a4) |- (a5);
  \draw[flecha] (a3) -- node[etiqueta,right]{no} (a5);
  \draw[flecha] (a5) -- (a6);
  \draw[flecha] (a6) -| node[etiqueta,above,pos=0.2]{sí} (a7);
  \draw[flecha] (a6) -| node[etiqueta,above,pos=0.2]{no} (a8);
  \draw[flecha] (a7) -- (a7 |- a9.north);
  \draw[flecha] (a8) -- (a8 |- a9.north);
  \draw[flecha] (a9) -- (a10);
  \draw[flecha] (a10) -- (a11);
  % salida "no" de la FIFO: pasa por fuera, a la izquierda
  \draw[flecha] (a2.west) -- node[etiqueta,above,pos=0.3]{no} (-4.4,0 |- a2.west)
                |- (a11.west);
\end{tikzpicture}
\caption{Diagrama de flujo de la rutina de fin de conversión del ADC. El
\emph{bypass} también se usa mientras \texttt{pte\_aux} es \texttt{NULL}
(antes de \texttt{init\_filtros()}).}
\label{fig:flujo-adc}
\end{figure}

\begin{figure}[H]
\centering
\begin{tikzpicture}[node distance=0.55cm and 0.4cm]
  % ---------------- ISR de SW3 ----------------
  \node[inicio] (b0) {\texttt{GPIO0\_INT\_0} (SW3)};
  \node[proceso, below=of b0] (b1) {\texttt{state\_counter} $\leftarrow$\\(\texttt{state\_counter}+1) mod 5};
  \node[proceso, below=of b1] (b2) {\texttt{setup\_new\_match}\\(\texttt{sample\_vel[state]})};
  \node[proceso, below=of b2] (b3) {\texttt{set\_state(state)}\\(color del LED)};
  \node[proceso, below=of b3] (b4) {\texttt{cambio\_filtro} $\leftarrow$ 1};
  \node[proceso, below=of b4] (b5) {Limpiar bandera};
  \draw[flecha] (b0) -- (b1);
  \draw[flecha] (b1) -- (b2);
  \draw[flecha] (b2) -- (b3);
  \draw[flecha] (b3) -- (b4);
  \draw[flecha] (b4) -- (b5);

  % ---------------- ISR de SW2 ----------------
  \begin{scope}[xshift=6.5cm]
    \node[inicio] (c0) {\texttt{GPIO0\_INT\_1} (SW2)};
    \node[proceso, below=of c0] (c1) {\texttt{filter\_counter} $\leftarrow$\\(\texttt{filter\_counter}+1) mod 5};
    \node[proceso, below=of c1] (c2) {\texttt{cambio\_filtro} $\leftarrow$ 1};
    \node[proceso, below=of c2] (c3) {Limpiar bandera};
    \draw[flecha] (c0) -- (c1);
    \draw[flecha] (c1) -- (c2);
    \draw[flecha] (c2) -- (c3);
  \end{scope}
\end{tikzpicture}
\caption{Diagramas de flujo de las rutinas de las teclas. Izquierda: cambio
de velocidad (SW3). Derecha: cambio de filtro (SW2).}
\label{fig:flujo-teclas}
\end{figure}

\subsection{Restricción de tiempo real}\label{sec:carga}

Todo el procesamiento de una muestra tiene que terminar antes de que llegue
la siguiente. Con el núcleo a $f_{\text{CPU}} = \qty{150}{\mega\hertz}$, los
ciclos disponibles por muestra son $f_{\text{CPU}}/f_s$, y el filtro consume
aproximadamente $c\,N$ ciclos, donde $c$ es el costo por coeficiente de
\texttt{arm\_fir\_q15()}. La fracción del período de muestreo que ocupa el
filtro es entonces
\begin{equation}
  U = \frac{c\, N\, f_s}{f_{\text{CPU}}},
  \label{eq:carga}
\end{equation}
y el sistema funciona en tiempo real sólo si $U < 1$ (con algo de margen para
el resto de la rutina y las teclas). La tabla~\ref{tab:carga} aplica la
ecuación~\eqref{eq:carga} a los filtros de la tabla~\ref{tab:taps} con
$c = \ciclosportap$ ciclos por coeficiente, la estimación que usa
\archivo{generar_filtros.m}.

\begin{table}[H]
\centering
\caption{Carga estimada de CPU, $U$ [\%], con $c = \ciclosportap$ ciclos por
coeficiente. En rojo, las combinaciones que no entran en tiempo real.}
\label{tab:carga}
\begin{tabular}{@{}l r r r r r@{}}
\toprule
\textbf{$f_s$} & \textbf{Ciclos/muestra} & \textbf{PB} & \textbf{PA} &
\textbf{PBanda} & \textbf{EB}\\
\midrule
\filacarga{8000} {64} {822} {822} {1270}
\filacarga{16000}{124}{1642}{1640}{2536}
\filacarga{22000}{170}{2256}{2254}{3486}
\filacarga{44000}{340}{4508}{4508}{6970}
\filacarga{48000}{370}{4918}{4916}{7604}
\bottomrule
\end{tabular}
\end{table}

Según esta estimación, a \qty{8}{\kilo\sps} los cuatro filtros entran con
holgura y el PB entra a todas las velocidades, pero el EB supera el período de muestreo a partir de
\qty{16}{\kilo\sps} y el PA y el PBanda a partir de \qty{22}{\kilo\sps}. La carga crece con $f_s^2$: al duplicar $f_s$ se duplica
la cantidad de coeficientes (sección~\ref{sec:taps}) y también la cantidad de
muestras por segundo.

Cuando $U > 1$ la rutina del ADC todavía no terminó cuando llega el disparo
siguiente. El disparo lo genera el hardware, así que el ADC sigue
convirtiendo, pero la interrupción queda pendiente y se atiende recién al
salir de la actual: se pierden muestras y la salida se actualiza a una tasa
menor que $f_s$. \completar{Describir lo que se observó en la placa en
esos casos (por ejemplo, salida distorsionada o LED/teclas que dejan de
responder).}

\begin{nota}[El costo por coeficiente depende de la compilación]
El valor $c = \ciclosportap$ es una estimación. El valor real depende del
nivel de optimización (la configuración \emph{Debug} de MCUXpresso compila
por defecto con \texttt{-O0}, también la biblioteca CMSIS-DSP), de los
estados de espera de la flash donde están los coeficientes y de la sobrecarga
de llamar a la función con una sola muestra. En la
sección~\ref{sec:verificacion} se mide; si difiere, basta cambiar
\texttt{\textbackslash ciclosportap} en el preámbulo de \archivo{main.tex} y la
tabla~\ref{tab:carga} se recalcula.
\end{nota}

%%%%%%%%%%%%%%%%%%%%%%%%%%%%%%%%%%%%%%%%%%%%%%%%%%%%%%%%%%%%%%%%%%%%%%%%%%%%%%%
\ejercicio{Implementación}\label{sec:implementacion}
%%%%%%%%%%%%%%%%%%%%%%%%%%%%%%%%%%%%%%%%%%%%%%%%%%%%%%%%%%%%%%%%%%%%%%%%%%%%%%%

La lógica propia está en \archivo{source/TP2.c} y los coeficientes en
\archivo{source/filtros.h}; el resto del proyecto es código generado por
Config Tools y drivers del SDK. El proyecto completo está referenciado en el
Anexo~\ref{sec:codigo}.

\subsection{Constantes y variables globales}

Los valores de match se mantienen en \texttt{sample\_vel}, igual que en el
Laboratorio 1. \texttt{filtros.h} define \texttt{NUM\_FS} (5),
\texttt{NUM\_TIPOS} (4) y \texttt{MAX\_TAPS}, que dimensionan la matriz de
instancias y el buffer de estado. Las variables que comparten las rutinas de
interrupción (\texttt{state\_counter}, \texttt{filter\_counter},
\texttt{cambio\_filtro} y el puntero \texttt{pte\_aux}) se declaran
\texttt{volatile}. Los buffers \texttt{buffer\_in} y \texttt{buffer\_out}
guardan las últimas 512 muestras de entrada y de salida, lo que permite
inspeccionarlas con el depurador.

\subsection{Programa principal}

\texttt{main()} inicializa pines y relojes, luego el DAC
(\texttt{init\_dac()}) y recién después los periféricos generados por Config
Tools, por el motivo explicado en la sección~\ref{sec:cambios-perif}. A
continuación inicializa las 20 instancias de filtro, fija el color del estado
inicial, activa \texttt{cambio\_filtro} para que la primera muestra aplique
la selección inicial y queda en un lazo vacío.

\subsection{Inicialización de los filtros}

\texttt{init\_filtros()} recorre la matriz de instancias llamando a
\texttt{arm\_fir\_init\_q15()} con los coeficientes y la cantidad de
\texttt{coef\_tabla} y \texttt{taps\_tabla}, el buffer de estado compartido y
\texttt{BLOCK\_SIZE}. Si alguna inicialización falla (por ejemplo, por una
cantidad impar de coeficientes), lo informa por la consola de depuración.

\subsection{Rutina del ADC}

El código~\ref{lst:isr} muestra el núcleo de la rutina del ADC. Antes de
procesar la muestra se aplica, si corresponde, el cambio de filtro pendiente;
después se convierte a q15, se filtra (o se copia) y se envía al DAC.

\begin{lstlisting}[caption={Procesamiento de una muestra en \texttt{ADC0\_IRQHANDLER()} (\texttt{source/TP2.c}).},label={lst:isr}]
if (LPADC_GetConvResult(ADC0_PERIPHERAL, &result, 0)) {
  if (cambio_filtro) {
    aplicar_seleccion();
  }

  /* ADC de 16 bits: 0..65535 -> Q15 centrado en 0 */
  buffer_in[index_buffer] = (q15_t)((int32_t)result.convValue - 32768);

  if (filter_counter == BYPASS || pte_aux == NULL) {
    buffer_out[index_buffer] = buffer_in[index_buffer];      /* bypass */
  } else {
    arm_fir_q15(pte_aux, &buffer_in[index_buffer],
                &buffer_out[index_buffer], BLOCK_SIZE);
  }

  /* Q15 -> 12 bits para el DAC */
  int32_t out_dac = ((int32_t)buffer_out[index_buffer] / 16) + ADC_MID;
  if (out_dac < 0)    out_dac = 0;
  if (out_dac > 4095) out_dac = 4095;
  DAC_SetData(DAC0, (uint16_t)out_dac);

  index_buffer = (index_buffer + 1) % BUFFER_SIZE;
}
\end{lstlisting}

\subsection{Rutinas de las teclas y cambio de velocidad}

La rutina de SW3 avanza \texttt{state\_counter} en forma circular,
reprograma el CTIMER0 con \texttt{setup\_new\_match()} y actualiza el LED con
\texttt{set\_state()}, igual que en el Laboratorio 1, y además activa
\texttt{cambio\_filtro} para que se carguen los coeficientes de la nueva
$f_s$. La rutina de SW2 avanza \texttt{filter\_counter} módulo 5 y activa la
misma bandera.

La macro \texttt{FILA\_FILTRO()} permite, poniendo \texttt{SOLO\_8K} en 1,
usar siempre los coeficientes diseñados para \qty{8}{\kilo\sps} sin importar
la velocidad activa. Sirve para observar cómo se desplazan en frecuencia las
bandas de un filtro digital cuando se lo usa a una $f_s$ distinta de la de
diseño: todas las frecuencias características escalan en proporción a $f_s$.

%%%%%%%%%%%%%%%%%%%%%%%%%%%%%%%%%%%%%%%%%%%%%%%%%%%%%%%%%%%%%%%%%%%%%%%%%%%%%%%
\ejercicio{Verificación}\label{sec:verificacion}
%%%%%%%%%%%%%%%%%%%%%%%%%%%%%%%%%%%%%%%%%%%%%%%%%%%%%%%%%%%%%%%%%%%%%%%%%%%%%%%

\subsection{Procedimiento}

\begin{enumerate}[nosep]
  \item Conectar el generador a \texttt{J4[2]} (ADC0\_A0) con una senoidal
        de \completar{amplitud} y un \emph{offset} que la mantenga entre 0 y
        \qty{3.3}{\volt}.
  \item Canal 1 del osciloscopio en la entrada, canal 2 en \texttt{J1[4]}
        (DAC0\_OUT).
  \item Con $f_s = \qty{8}{\kilo\sps}$, recorrer los cinco estados de SW2 y,
        para cada filtro, barrer la frecuencia del generador midiendo la
        amplitud de la salida. Calcular la ganancia
        $|H| = 20\log_{10}(V_{\text{out}}/V_{\text{in}})$ y compararla con la
        respuesta diseñada.
  \item Repetir para otras $f_s$ y verificar que la respuesta en hertz se
        mantiene (o, si no entra en tiempo real, documentar qué ocurre).
  \item Medir el tiempo de ejecución de la rutina del ADC
        \completar{método: por ejemplo, conmutar un pin GPIO al entrar y salir
        de la rutina y medir el pulso con el osciloscopio, o leer el contador
        de ciclos \texttt{DWT->CYCCNT}}.
\end{enumerate}

\subsection{Respuesta en frecuencia medida}

Las frecuencias de prueba se eligen alrededor de los bordes de la consigna:
\qty{50}{\hertz} (centro del EB), unos \qty{20}{\hertz} (rechazo del PA y del
PBanda), \qty{35}{\hertz} (borde del PA y del PBanda), \qty{1}{\kilo\hertz}
(banda de paso de todos los filtros salvo el EB), y \qty{3.5}{\kilo\hertz} y
\qty{3.9}{\kilo\hertz} (bordes del PBanda y del PB). Para las frecuencias de
pocas decenas de hertz hay que usar acoplamiento DC en el osciloscopio: el
acoplamiento AC tiene un corte de algunos hertz que falsearía la medición.
\completar{Indicar en el texto o en cada captura la frecuencia de la señal de
entrada usada.}

\begin{table}[H]
\centering
\caption{Ganancia medida a $f_s = \qty{8}{\kilo\sps}$.}
\label{tab:respuesta}
\begin{tabular}{@{}r r r r r@{}}
\toprule
\textbf{$f$ [Hz]} & \textbf{PB [dB]} & \textbf{PA [dB]} &
\textbf{PBanda [dB]} & \textbf{EB [dB]}\\
\midrule
\num{20}   & \completar{} & \completar{} & \completar{} & \completar{}\\
\num{35}   & \completar{} & \completar{} & \completar{} & \completar{}\\
\num{50}   & \completar{} & \completar{} & \completar{} & \completar{}\\
\num{1000} & \completar{} & \completar{} & \completar{} & \completar{}\\
\num{3500} & \completar{} & \completar{} & \completar{} & \completar{}\\
\num{3900} & \completar{} & \completar{} & \completar{} & \completar{}\\
\bottomrule
\end{tabular}
\end{table}

\captura{img/osc-bypass.jpeg}{Osciloscopio}{Entrada (canal 1) y salida del
DAC (canal 2) en \emph{bypass}.}{fig:osc-bypass}

\captura{img/osc-pb.jpeg}{Osciloscopio}{Entrada y salida con el filtro pasa
bajos.}{fig:osc-pb}

\captura{img/osc-pa.jpeg}{Osciloscopio}{Entrada y salida con el filtro pasa
altos.}{fig:osc-pa}

\captura{img/osc-pbanda.jpeg}{Osciloscopio}{Entrada y salida con el filtro
pasa banda.}{fig:osc-pbanda}

\captura{img/osc-eb.jpeg}{Osciloscopio}{Entrada y salida con el filtro
elimina banda.}{fig:osc-eb}

\subsection{Tiempo de ejecución}

\begin{table}[H]
\centering
\caption{Duración medida de la rutina del ADC.}
\label{tab:tiempo}
\begin{tabular}{@{}l l r r r@{}}
\toprule
\textbf{$f_s$} & \textbf{Filtro} & \textbf{$N$} & \textbf{Duración medida} &
\textbf{Ciclos por coeficiente}\\
\midrule
\qty{8}{\kilo\sps} & Bypass & --       & \completar{} & --\\
\qty{8}{\kilo\sps} & PB     & \num{64}   & \completar{} & \completar{}\\
\qty{8}{\kilo\sps} & EB     & \num{1270} & \completar{} & \completar{}\\
\bottomrule
\end{tabular}
\end{table}

\captura{img/osc-tiempo-isr.jpeg}{Osciloscopio}{Pulso de un pin de prueba
mientras se ejecuta la rutina del ADC.}{fig:osc-tiempo}

\subsection{Cumplimiento de los requerimientos}

\begin{table}[H]
\centering
\caption{Verificación de los requerimientos de la consigna.}
\label{tab:cumplimiento}
\begin{tabular}{@{}l p{6.8cm} l l@{}}
\toprule
\textbf{Req.} & \textbf{Cómo se implementa} & \textbf{Evidencia} & \textbf{Resultado}\\
\midrule
\ref{req:fs}           & Cadena CTIMER0 + INPUTMUX + LPADC0 del Laboratorio~1;
                         SW3 elige la $f_s$; \texttt{buffer\_in[512]}
                       & Tabla~\ref{tab:fs} & \completar{}\\
\ref{req:diseno}       & 20 filtros (4 tipos $\times$ 5 $f_s$) diseñados en
                         \texttt{fdatool} y exportados con
                         \archivo{generar_filtros.m}
                       & Tabla~\ref{tab:specs} & \completar{}\\
\ref{req:cmsis}        & \texttt{arm\_fir\_q15()} con \texttt{BLOCK\_SIZE = 1},
                         de \texttt{buffer\_in} a \texttt{buffer\_out}
                       & Tabla~\ref{tab:respuesta} & \completar{}\\
\ref{req:tecla-filtro} & SW2 recorre \emph{bypass} y F1--F4
                         (\texttt{filter\_counter} módulo 5)
                       & Demostración & \completar{}\\
\ref{req:dac}          & \texttt{DAC\_SetData()} con la muestra de
                         \texttt{buffer\_out} en cada interrupción del ADC
                       & Figura~\ref{fig:osc-bypass} & \completar{}\\
\ref{req:osc}          & Entrada y salida en los canales 1 y 2 del osciloscopio
                       & Figuras~\ref{fig:osc-pb} a~\ref{fig:osc-eb} & \completar{}\\
\bottomrule
\end{tabular}
\end{table}

\ejercicio{Conclusiones}

Se implementó sobre la FRDM-MCXN947 una aplicación que filtra en tiempo real
la señal adquirida por el LPADC0 con filtros FIR pasa bajos, pasa altos, pasa
banda y elimina banda, diseñados en MATLAB y ejecutados con
\texttt{arm\_fir\_q15()} de CMSIS-DSP en aritmética q15. Con SW3 se recorren
las cinco frecuencias de muestreo del Laboratorio 1 y con SW2 se elige entre
la señal sin filtrar y cada uno de los cuatro filtros; para cada combinación
se usan coeficientes diseñados para la $f_s$ activa.
\completar{Resumir el resultado de la tabla~\ref{tab:cumplimiento}.}

Las decisiones principales del diseño fueron:
\begin{itemize}[nosep]
  \item \textbf{Un juego de coeficientes por $f_s$}, para que la respuesta en
        frecuencia en hertz no dependa de la velocidad elegida. El precio es
        que la cantidad de coeficientes, y con ella la carga de CPU, crecen
        con $f_s$.
  \item \textbf{ADC en \qty{16}{\bit}}, cuyo resultado, centrado, es
        directamente un \texttt{q15\_t}.
  \item \textbf{Un único buffer de estado compartido} y un puntero al filtro
        activo, en lugar de 20 buffers de estado.
  \item \textbf{Cambio de filtro diferido a la rutina del ADC}, con borrado
        del estado, para que nunca se cambie de filtro en medio de un cálculo.
  \item \textbf{Diseño por ventana de Kaiser}, que es robusto para filtros de
        miles de coeficientes, donde el método de Parks-McClellan puede no
        converger. A cambio, Kaiser requiere entre un \qty{30}{\percent} y un
        \qty{35}{\percent} más de coeficientes que un diseño equiripple con las
        mismas especificaciones.
\end{itemize}

El método de diseño no es, sin embargo, la causa de fondo de la carga de CPU.
Según las ecuaciones~\eqref{eq:kaiser} y~\eqref{eq:equiripple}, con cualquier
método el orden de un FIR es proporcional a $f_s/\Delta f$. Las
especificaciones del PA, del PBanda y del EB tienen transiciones de
\qtyrange{7.5}{15}{\hertz}, y con $f_s$ de decenas de kilohertz eso exige
miles de coeficientes por muestra. Pasar a un diseño equiripple reduciría la
carga en alrededor de un tercio, pero no alcanza para que esos filtros entren
en tiempo real a partir de \qty{22}{\kilo\sps}: el problema requiere un factor
de 5 a 10, no de un tercio. Las mejoras con mayor impacto, dentro del
procesamiento por muestras que pide la consigna, serían:
\begin{itemize}[nosep]
  \item \textbf{Procesamiento multitasa}: las bandas exigentes (\qty{35}{\hertz}
        y \qty{50}{\hertz}) están muy por debajo de $f_s/2$, así que se pueden
        filtrar a una tasa reducida, después de diezmar, y volver a interpolar.
        Para el PA y el EB también se puede usar la estructura complementaria
        $y = x - h_{\text{PB}} * x$, con el pasa bajos angosto calculado a baja
        tasa.
  \item \textbf{Aprovechar la simetría de los coeficientes}: sumar primero las
        muestras que comparten coeficiente, $x[n-k] + x[n-N+1+k]$, reduce las
        multiplicaciones a la mitad. CMSIS-DSP no ofrece esta variante en q15,
        así que habría que implementarla.
  \item \textbf{Revisar los bordes elegidos}: la consigna sólo fija
        $f_c = \qty{35}{\hertz}$ para el PA. Bajar $f_{\text{stop}}$ de
        \qty{20}{\hertz} a \qty{10}{\hertz} ensancharía la transición de
        \qty{15}{\hertz} a \qty{25}{\hertz} y reduciría el orden en alrededor
        de un \qty{40}{\percent}.
  \item \textbf{Compilar con optimización} (\texttt{-O2} o \texttt{-O3}) para
        bajar los ciclos por coeficiente.
\end{itemize}

\completar{Comparar la respuesta medida (tabla~\ref{tab:respuesta}) con la
diseñada (figuras~\ref{fig:fda-pb} a~\ref{fig:fda-eb}) y comentar el
retardo observado entre entrada y salida frente al de la
ecuación~\eqref{eq:retardo}.}

\completar{Discutir la restricción de tiempo real con los valores medidos de
la tabla~\ref{tab:tiempo}: qué combinaciones de $f_s$ y filtro funcionan y
cuáles no, y posibles mejoras (por ejemplo, relajar las especificaciones para
bajar el orden, compilar con optimización, procesar por bloques o usar
\texttt{arm\_fir\_fast\_q15()}).}

%%%%%%%%%%%%%%%%%%%%%%%%%%%%%%%%%%%%%%%%%%%%%%%%%%%%%%%%%%%%%%%%%%%%%%%%%%%%%%%
\clearpage
\appendix
\section{Código fuente}\label{sec:codigo}
%%%%%%%%%%%%%%%%%%%%%%%%%%%%%%%%%%%%%%%%%%%%%%%%%%%%%%%%%%%%%%%%%%%%%%%%%%%%%%%

El proyecto completo de MCUXpresso (incluidos los archivos generados por
Config Tools, los drivers del SDK, CMSIS-DSP y \archivo{source/filtros.h})
está en la carpeta \archivo{TP2/} del repositorio
\url{https://github.com/Mati-Costamagna/procesamiento-digital-de-senales}.
\completar{Agregar \archivo{generar_filtros.m} al repositorio e indicar su
ubicación.}

\end{document}
